\documentclass[11pt]{amsart}
\usepackage[utf8]{inputenc}
\usepackage[T1]{fontenc}
\usepackage[french]{babel}
\usepackage{amsmath,amssymb,amsthm,amscd}
\usepackage{mathrsfs}
\usepackage{enumerate}

\theoremstyle{plain}
\newtheorem{theorem}{Th\'eor\`eme}[section]
\newtheorem{lemma}[theorem]{Lemme}
\newtheorem{proposition}[theorem]{Proposition}
\newtheorem{corollary}[theorem]{Corollaire}

\theoremstyle{definition}
\newtheorem{definition}[theorem]{D\'efinition}
\newtheorem{remark}[theorem]{Remarque}
\newtheorem{remarks}[theorem]{Remarques}
\newtheorem{example}[theorem]{Exemple}

\DeclareMathOperator{\Gal}{Gal}
\DeclareMathOperator{\Ind}{Ind}
\DeclareMathOperator{\Res}{Res}
\DeclareMathOperator{\Hom}{Hom}
\DeclareMathOperator{\End}{End}
\DeclareMathOperator{\Aut}{Aut}
\DeclareMathOperator{\Spec}{Spec}
\DeclareMathOperator{\Ker}{Ker}
\DeclareMathOperator{\GL}{GL}
\DeclareMathOperator{\Id}{Id}
\DeclareMathOperator{\Lie}{Lie}
\DeclareMathOperator{\Tr}{Tr}
\DeclareMathOperator{\rank}{rank}
\DeclareMathOperator{\sw}{sw}
\DeclareMathOperator{\gr}{Gr}

\begin{document}
\title{Les constantes des \'equations fonctionnelles des fonctions $L$}
\author{P. Deligne}
\maketitle

\setcounter{section}{4}

\section{Formulaire}
\setcounter{theorem}{10}

\begin{remark}[5.11]
Soit $K$ un corps global. Pour $V$ une repr\'esentation du groupe de Weil global de $K$, $\psi$ un caract\`ere additif non trivial de $\mathbb{A}/K$ et $dx = \otimes dx_v$ la mesure de Tamagawa sur $\mathbb{A}_K$, on pose
\begin{equation}
L(V) = \prod_v L(V_v)
\tag{5.11.0}
\end{equation}
\begin{equation}
\varepsilon(V) = \prod_v \varepsilon(V_v, \psi_v, dx_v)
\tag{5.11.1}
\end{equation}
L'\'equation fonctionnelle globale s'\'ecrit
\begin{equation}
L(V) = \varepsilon(V) \cdot L(V^*).
\tag{5.11.2}
\end{equation}
\end{remark}

\noindent\textbf{D\'emonstrations.} 5.2, 5.3, 5.4 et 5.5 r\'esultent de 4.1 (1) et (2), (3), (4) (et 3.3). Il suffit de prouver (5.7.1) pour $V$ d\'efini par un caract\`ere; la formule r\'esulte alors de (5.8.1) et de la formule d'inversion de Fourier. (5.7.2) r\'esulte de (5.7.1) et de (5.3), (5.4). 5.9 r\'esulte de (3.4.3.1) et (5.10) de (5.8.2). Enfin, 5.11 r\'esulte de ce qui pr\'ec\`ede et de (3.11.4) par des arguments connus [1], [19].

La restriction de 4.1 au cas des extensions et des caract\`eres mod\'er\'ement ramifi\'es \'equivaut aux identit\'es suivantes sur les sommes de Gauss (5.10.2). On y d\'esigne par $k$ un corps fini \`a $q$ \'el\'ements et par $\psi$ un caract\`ere additif non trivial de $k$.

\begin{theorem}[5.12 --- Hasse-Davenport]
Soient $k'$ une extension de $k$ de degr\'e $n$ et $\chi$ un caract\`ere de $k^*$. On a
\[
\tau(\chi \circ N_{k'/k}, \psi \circ \operatorname{Tr}_{k'/k}) = (-1)^{n-1} \, \tau(\chi, \psi)^n.
\]
\end{theorem}

\begin{theorem}[5.13]
Soient $\chi$ un caract\`ere de $k^*$ et $n$ un entier premier \`a $q$. Soit $X$ un ensemble de repr\'esentants des classes d'isomorphie de couples $(k', \chi')$ form\'es d'une extension $k'$ de $k$ et d'un caract\`ere $\chi'$ de $k'^*$ tel que
\begin{enumerate}
\item[(a)] $k'$ est interm\'ediaire entre $k$ et $k'$, i.e. $\chi' \circ N_{k'/k''}$ pour $k \subset k'' \subset k'$, $0 < [k'':k] < [k':k]$.
\item[(b)] $\chi' \circ N_{k'/k''} \not= \chi'$ pour $0 < [k'':k] < [k':k]$.
\end{enumerate}
Alors,
\begin{equation}
\prod_{(k', \chi') \in X} \tau(\chi', \psi \circ \operatorname{Tr}_{k'/k}) = \tau(\chi, \psi).
\tag{5.13.1}
\end{equation}
\end{theorem}

Cette identit\'e se divise en les deux suivantes.

\begin{theorem}[5.14 --- Langlands {[9]} 7.8]
Soient $\ell$ un nombre premier \`a $q$, $f$ l'ordre de $q$ dans $(\mathbb{Z}/\ell\mathbb{Z})^*$, $k'$ une extension de $k$ de degr\'e $f$, $T$ un ensemble de repr\'esentants des orbites de $\Gal(k'/k)$ dans les caract\`eres (non triviaux) d'ordre $\ell$ de $k'^*$ et $\chi$ un caract\`ere de $k^*$. On a
\[
\prod_{\theta \in T} \tau(\theta, \psi \circ \operatorname{Tr}_{k'/k}) \cdot \tau(\chi, \psi) = \tau(\chi \circ N_{k'/k}, \psi \circ \operatorname{Tr}_{k'/k}) \cdot \prod_{\theta \in T} \tau(\chi \circ N_{k'/k} \cdot \theta, \psi \circ \operatorname{Tr}_{k'/k}).
\]
\end{theorem}

\begin{theorem}[5.15 --- Langlands {[9]} 7.9]
Soient $\ell$ un nombre premier divisant $q-1$, $T$ l'ensemble des caract\`eres d'ordre $\ell$ de $k^*$, $\chi$ un caract\`ere de $k^*$ qui n'est pas une puissance $\ell$-i\`eme, $k'$ une extension de degr\'e $\ell$ de $k$ et $\chi'$ un caract\`ere de $k'^*$ tel que $\chi'(x^\ell) = \chi \circ N_{k'/k}(x)$. On a
\[
\prod_{\theta \in T} \tau(\theta, \psi) \cdot \tau(\chi, \psi) = \tau(\chi', \psi \circ \operatorname{Tr}_{k'/k}).
\]
\end{theorem}

\begin{proof}[5.16 --- D\'emonstrations]
Si $K'/K$ est une extension mod\'er\'ement ramifi\'ee de corps locaux non archim\'ediens et $\psi$ un caract\`ere additif de $K$ tel que $n(\psi) = -1$, on a encore $n(\psi \circ \operatorname{Tr}_{K'/K}) = -1$, de sorte que 5.10 s'applique tant \`a $K$ qu'\`a $K'$.

Soit $V$ une repr\'esentation mod\'er\'ement ramifi\'ee de $W(\overline{K}/K)$. On peut toujours l'\'ecrire comme somme de repr\'esentations induites par des caract\`eres mod\'er\'es $\chi_i$ d'extensions non ramifi\'es $K_i/K$:
\[
V = \dim(V) \cdot [1] + \sum_i \Ind([\chi_i] - [1]) + \sum_i (\Ind([1]) - [K_i:K][1]).
\]
Soient $k_i$ le corps r\'esiduel de $K_i$ et $\overline{\chi}_i$ le caract\`ere de $k_i^*$ d\'efini par $\chi_i$. Pour $dx$ une mesure de Haar sur $K$ telle que $\int_{\mathcal{O}} dx = 1$, on trouve
\begin{equation}
\varepsilon_0(V,\psi,dx) = (-1)^{\dim V} \prod_i \tau(\overline{\chi}_i, \psi_k \circ \operatorname{Tr}_{k_i/k}).
\tag{5.16.1}
\end{equation}
La formule (5.12) s'obtient en prenant une extension non ramifi\'ee $K'/K$, et en \'ecrivant que, pour $\chi$ un caract\`ere (ici mod\'er\'ement ramifi\'e) de $K^*$, on a $\sum_{\mu} [\mu\chi] - [\mu]$ o\`u $\mu$ parcourt les caract\`eres non ramifi\'es d'ordre divisant $[K':K]$. Cette formule (5.12) assure que le second membre de (5.16.1) ne d\'epend que de $V$.

Les identit\'es 5.13 s'obtiennent en prenant une extension totalement ramifi\'ee de degr\'e $n$ $K'/K$, et en conjuguant 5.6.1 et (5.16.1) pour $V$ induite par un caract\`ere mod\'er\'e $\chi$ de $K'$.

(5.14) et (5.15) sont le cas particulier de (5.13.1) o\`u $n$ est un nombre premier $\ell$. Dans 5.14 (resp. 5.15), on suppose que le caract\`ere $\chi$ de $K'^*$ est (resp. n'est pas) une puissance $\ell$-i\`eme.
\end{proof}

\section{R\'eduction modulo $\ell$ des constantes locales}

Soient $K$ un corps local non archim\'edien, $\overline{K}$ une cl\^oture s\'eparable de $K$ et reprenons les notations 2.2. Soit $\Lambda$ un anneau commutatif dans lequel $p$ soit inversible. On simplifierait l'expos\'e, sans perte de g\'en\'eralit\'e, en supposant que $\Lambda$ est local, voire un corps. Pour $x \in K$, $\|x\|$ a par hypoth\`ese un sens dans $\Lambda$.

\begin{definition}[6.1]
Une mesure de Haar $\mu$ sur $K$, \`a valeurs dans un $\mathbb{Z}[1/p]$-module $M$ est une fonction simplement additive de parties compactes ouvertes de $K$, \`a valeurs dans $M$, et invariante par translation. La fonction additive d\'efinie par $\mu_0(a + \pi^n\mathcal{O}) = \|\pi\|^n$ est l'unique mesure de Haar \`a valeurs dans $\mathbb{Z}[1/p]$ telle que $\mu_0(\mathcal{O}) = 1$, et toute mesure de Haar \`a valeurs dans $M$ est de la forme $m \cdot \mu_0$ ($m = \mu_0(\mathcal{O}) \in M$). On peut int\'egrer une fonction localement constante \`a support compact \`a valeurs dans $\Lambda$ contre une mesure de Haar \`a valeurs dans $\Lambda$.
\end{definition}

\begin{definition}[6.2]
Soient $\rho: W(\overline{K}/K) \to \GL_\Lambda(V)$ une repr\'esentation du groupe de Weil sur un $\Lambda$-module libre (ou projectif) $V$, et $J$ un sous-groupe distingu\'e ouvert de $I$ sur lequel $\rho$ soit trivial. Pour tout nombre premier $\ell \not= p$, on sait que la repr\'esentation de Swan $\sw_{I/J}$ peut se r\'ealiser comme un $\mathbb{Z}_\ell[I/J]$-module projectif. Si $\Lambda$ est une $\mathbb{Z}_\ell$-alg\`ebre, $\Hom_{I/J}(\sw_{I/J}, V)$ est donc un $\Lambda$-module projectif. Si $\Lambda$ est de plus local, il est libre; son rang est le conducteur de Swan de $V$. Plus g\'en\'eralement, on v\'erifie qu'on peut toujours d\'efinir le conducteur de Swan de $V$ comme une fonction localement constante \`a valeurs enti\`eres sur $\Spec(\Lambda)$. Ce conducteur est invariant par extension de l'anneau des scalaires $\Lambda$.
\end{definition}

\begin{definition}[6.3]
Si $\zeta$ est une racine $(p^n)$-i\`eme de l'unit\'e dans $\Lambda$, $\zeta$ est une fonction localement constante sur $\Spec(\Lambda)$: deux quelconques des polyn\^omes cyclotomiques $\Phi_{p^i}(X)$ ($1 \leq i \leq n$) engendrent l'id\'eal trivial de $\Lambda[X]$ (puisque $p \in \Lambda^*$), donc $\Lambda$ est produits d'anneaux $\Lambda_i$ ($1 \leq i \leq n$), la $i$\`eme coordonn\'ee de $\zeta$ v\'erifiant $\Phi_{p^i}(\zeta) = 0$. Sur $\Lambda_i$, $\zeta$ est d'ordre $p^i$.

Si $\psi$ est un caract\`ere additif (continu) de $K$ \`a valeurs dans $\Lambda^*$, il est \`a valeurs dans les racines $(p^n)$-i\`emes de l'unit\'e. Ceci permet de d\'efinir $n(\psi)$ comme une fonction localement constante sur $\Spec(\Lambda)$; en un id\'eal premier $\mathfrak{p}$, $n(\psi)$ est $+\infty$ si $\psi$ est trivial mod $\mathfrak{p}$, et le plus grand entier $n$ tel que $\psi|_{\pi^{-n}\mathcal{O}}$ soit trivial mod $\mathfrak{p}$ (ou dans le localis\'e $\Lambda_{\mathfrak{p}}$, c'est pareil). Par abus de langage, on dit que $n(\psi)$ est non trivial si $n(\psi)$ n'est nulle part $+\infty$.
\end{definition}

\begin{definition}[6.4]
Pour $\psi$ un caract\`ere additif non trivial de $K$ \`a valeurs dans $\Lambda$, $\chi \in \Hom(K^*, \Lambda^*)$, et $dx$ une mesure de Haar \`a valeurs dans $\Lambda$, on peut maintenant d\'efinir $\varepsilon_0(\chi,\psi,dx)$ par la formule (3.4.3.4) (plus pr\'ecis\'ement, si $\Spec(\Lambda)$ n'est pas connexe, on commence par \'ecrire $\Lambda = \prod \Lambda_i$, avec $n(\psi)$ et $\sw(\chi)$ constant sur $\Spec(\Lambda_i)$. On d\'efinit alors la $i$\`eme projection de $\varepsilon_0$ par (3.4.3.4)).
\end{definition}

\begin{theorem}[6.5]
Il existe une et une seule fonction $\varepsilon_0$ du type suivant.
\begin{enumerate}
\item[(a)] Pour $K$ un corps local non archim\'edien, $\overline{K}$ une cl\^oture s\'eparable de $K$, $\Lambda$ un corps de caract\'eristique $\not= p$, $V$ une repr\'esentation de $W(\overline{K}/K)$ dans un vectoriel de dimension finie sur $\Lambda$, $\psi$ un caract\`ere additif non trivial \`a valeurs dans $\Lambda$ et $dx$ une mesure de Haar (non nulle) \`a valeurs dans $\Lambda$, on a $\varepsilon_0 \in \Lambda^*$.
\item[(b)] $\varepsilon_0$ est invariant par extension du corps $\Lambda$.
\item[(c)] Soit $\Lambda$ un anneau de valuation discr\`ete de caract\'eristique r\'esiduelle $\not= p$, $\eta$ son corps des fractions et $s = \Lambda/\mathfrak{m}$ son corps r\'esiduel. Pour $V$ une repr\'esentation de $W(\overline{K}/K)$ dans un $\Lambda$-module libre de rang fini, et pour $\psi$ et $dx$ \`a valeurs dans $\Lambda$, on a $\varepsilon_0(V,\psi,dx) \in \Lambda^*$, et $\overline{\varepsilon_0} \equiv \varepsilon_0(\overline{V},\overline{\psi},\overline{dx}) \pmod{\mathfrak{m}}$.
\item[(d)] Les conditions (1)--(3) de 4.1 sont v\'erifi\'ees.
\item[(e)] Pour $\dim V = 1$, $\varepsilon_0$ est donn\'e par 6.4.
\end{enumerate}
\end{theorem}

\begin{proof}
Par descente galoisienne, il suffit de traiter le cas o\`u $\Lambda$ est un corps s\'eparablement clos. Si $\Lambda$ est de caract\'eristique 0, l'\'enonc\'e, \'etant purement alg\'ebrique, r\'esulte du cas $\Lambda = \mathbb{C}$ trait\'e en 4.1. Supposons donc $\Lambda$ de caract\'eristique $\ell \not= 0$.

Soient $G$ un quotient fini de $\Gal(\overline{K}/K)$ et $\Lambda$ un anneau de valuation discr\`ete d'in\'egale caract\'eristique de corps r\'esiduel $\Lambda/\mathfrak{m} = \Lambda$. On suppose $\psi$ rel\`eve en un caract\`ere additif \`a valeurs dans $\Lambda^*$, $dx$ en une mesure de Haar \`a valeurs dans $\Lambda$, et que le corps des fractions $\eta$ de $\Lambda$ est assez gros (1.4) relativement \`a $G$.

Pour tout caract\`ere $\chi$ d'une sous-groupe de $G$ \`a valeurs dans $\eta^*$, en fait dans $\Lambda^*$, la constante $\varepsilon_0$ correspondante est clairement dans $\Lambda$. D'apr\`es 5.7.1, elle est m\^eme dans $\Lambda^*$. On en d\'eduit que, pour toute repr\'esentation de $G$ sur $\eta$, on a $\varepsilon_0(V_\eta,\psi,dx) \in \Lambda^*$.

On sait que l'homomorphisme de d\'ecomposition $d: R_\eta(G) \to R_\Lambda(G)$ est surjectif (les repr\'esentations virtuelles se rel\`event). Pour $\overline{V} \in R_\Lambda(G)$, et $V \in R_\eta(G)$ tel que $d(V) = \overline{V}$, on pose $\varepsilon_0(\overline{V},\overline{\psi},\overline{dx}) =$ r\'eduction mod $\mathfrak{m}$ de $\varepsilon_0(V,\psi,dx)$.

D'apr\`es 1.8, pour l\'egitimer cette d\'efinition, il suffit de v\'erifier que
\begin{flushleft}
(*) Pour $H \subset G$ d\'efinissant une extension $L$ de $K$, $\psi_L$ et $\chi', \chi''$ deux caract\`eres de $H$ \`a valeurs dans $\eta^*$, si $\chi' \equiv \chi'' \pmod{\mathfrak{m}}$, on a $\varepsilon_0(\chi',\psi_L,dx) \equiv \varepsilon_0(\chi'',\psi_L,dx) \pmod{\mathfrak{m}}$.
\end{flushleft}

Les conducteurs de Swan de $\chi'$ et $\chi''$ sont tous deux \'egaux au conducteur de Swan de $\chi'$ mod $\mathfrak{m}$. Pour $y = \pi^{1+\sw(\chi')+n(\psi_L)}$, on a donc
\[
\varepsilon_0 = \int_{y^{-1}\mathcal{O}_L^*} \chi'^{-1}(x) \, \psi_L(x) \, dx.
\]
Num\'erateur et d\'enominateur sont des unit\'es, et sont clairement congrus mod $\mathfrak{m}$. On a donc $\varepsilon_0 \equiv 1$.

Pour $V$ une repr\'esentation de $\Gal(\overline{K}/K)$, on obtient ainsi une fonction ayant les propri\'et\'es annonc\'ees. Ses propri\'et\'es formelles sont les m\^emes que celles explicit\'ees dans le formulaire 5.1 dans le cas $\Lambda = \mathbb{C}$. Pour passer de $\Gal(\overline{K}/K)$ \`a $W(\overline{K}/K)$ on peut donc reprendre l'argument utilis\'e en 4.9, \`a l'aide de la deuxi\`eme formule (5.5.1).
\end{proof}

\begin{remark}[6.6]
Dans tout ce num\'ero, il a \'et\'e essentiel de travailler avec $\varepsilon_0$ plut\^ot qu'avec $\varepsilon$. D'abord pour avoir \`a utiliser le conducteur de Swan (invariant par r\'eduction mod $\ell$) plut\^ot que le conducteur d'Artin, ensuite pour n'avoir pas \`a distinguer, dans la d\'efinition de $\varepsilon_0(\chi,\psi,dx)$, entre les cas non ramifi\'e et mod\'er\'ement ramifi\'e. La signification de $\varepsilon_0$ appara\^itra plus clairement au num\'ero suivant.
\end{remark}

\section{Fonctions $L$ modulaires}

\begin{definition}[7.1]
Soit $X$ une courbe projective non singuli\`ere sur $\mathbb{F}_p$. On suppose $X$ irr\'eductible, on note $K$ son corps de fonctions rationnelles et on identifie points ferm\'es de $X$ et places de $K$. On ne suppose pas $X$ g\'eom\'etriquement irr\'eductible, i.e. $\mathbb{F}_p$ alg\'ebriquement ferm\'e dans $K$.
\end{definition}

\begin{definition}[7.2]
Soit $\Lambda$ un anneau noeth\'erien. On sait ce qu'est un faisceau constructible de $\Lambda$-modules sur $X$ (SGA 4 IX 2.3). Pour tout $x \in X$, et pour $\overline{k(x)}$ une extension s\'eparablement close de $k(x)$, un tel faisceau $\mathcal{G}$ a une fibre $\mathcal{G}_{\overline{k(x)}}$ en $x$, qui est un $\Lambda$-module de type fini. Ce dernier ne d\'epend que de la cl\^oture s\'eparable de $k(x)$ dans $\overline{k(x)}$, et $\Gal(\overline{k(x)}/k(x))$ agit sur $\mathcal{G}_{\overline{k(x)}}$ par transport de structure.

De plus, pour $v$ une place de $K$ (point ferm\'e de $X$), et $\overline{K}_v$ une cl\^oture alg\'ebrique de $K_v$, on dispose d'une fl\`eche de sp\'ecialisation $\mathcal{G}_{\overline{k(v)}} \to \mathcal{G}_{\overline{K}_v}^{I_v}$. Cette fl\`eche est \'equivariante; son image est donc contenue dans $\mathcal{G}_{\overline{K}_v}^{I_v}$. On dit que $\mathcal{G}$ est non ramifi\'e en $v$ si $\mathrm{sp}_v$ est un isomorphisme.

Un faisceau constructible est presque partout non ramifi\'e. Une fois choisies une cl\^oture s\'eparable $\overline{K}$ de $K$ et une place de $\overline{K}$ au-dessus de chaque place de $K$, on peut r\'eciproquement d\'efinir un faisceau constructible de $\Lambda$-modules comme suit:
\begin{enumerate}
\item[(a)] On se donne les fibres $\mathcal{G}_{\overline{K}}$ et $\mathcal{G}_{\overline{k(v)}}$ (respectivement des $\Lambda$-$\Gal(\overline{K}/K)$- ou $\Lambda$-$\Gal(\overline{k(v)}/k(v))$-modules, de type fini sur $\Lambda$).
\item[(b)] On se donne les fl\`eches de sp\'ecialisation (\'equivariantes pour l'action des groupes de d\'ecomposition) $\mathcal{G}_{\overline{k(v)}} \to \mathcal{G}_{\overline{K}}^{I_v}$.
\end{enumerate}
Pour que (a), (b) d\'efinissent un faisceau, il suffit que les $\mathrm{sp}_v$ soient presque tous des isomorphismes.
\end{definition}

\begin{definition}[7.3]
Il nous sera commode de modifier la d\'efinition pr\'ec\'edente en rempla\c{c}ant les groupes de Galois par des groupes de Weil. Soit $Y$ un sch\'ema sur $\mathbb{F}_p$. Soient $\overline{\mathbb{F}}_p$ une cl\^oture alg\'ebrique de $\mathbb{F}_p$ et $\overline{Y} = Y \otimes_{\mathbb{F}_p} \overline{\mathbb{F}}_p$. Un faisceau de Weil (resp. un faisceau de Weil de $\Lambda$-modules constructible) $\mathcal{G}$ sur $Y$ est un faisceau (resp. un faisceau de $\Lambda$-modules constructible) $\mathcal{G}$ sur $\overline{Y}_{\text{\'et}}$, plus une action de Frobenius sur $\mathcal{G}$, au-dessus de son action sur $\overline{Y}$.

\emph{Remarque} (inutile): Les faisceaux de Weil en ensembles sur $Y$ forment un topos $Y_{\mathrm{wet}}$. On peut le d\'ecrire comme un 2-produit fibr\'e de topos: identifiant $\Spec(\overline{\mathbb{F}}_p)_{\text{\'et}}$ au topos classifiant $B_{\widehat{\mathbb{Z}}}$, on a $Y_{\mathrm{wet}} = Y_{\text{\'et}} \times_{B_{\widehat{\mathbb{Z}}}} B_{\mathbb{Z}}$.

Nous appellerons simplement faisceaux les faisceaux de Weil sur $\overline{Y}_{\text{\'et}}$; les faisceaux sur $Y_{\text{\'et}}$ seront appel\'es faisceaux \'etales.
\end{definition}

\begin{definition}[7.4]
Un faisceau 7.3 de $\Lambda$-modules constructible peut se d\'ecrire en termes galoisiens comme en 7.2, en rempla\c{c}ant partout $\Gal(\overline{K}/K)$, $\Gal(\overline{K}_v/K_v)$ et $\Gal(\overline{k(v)}/k(v))$ par $W(\overline{K}/K)$, $W(\overline{K}_v/K_v)$ et $W(\overline{k(v)}/k(v)) = \mathbb{Z}$.
\end{definition}

\begin{definition}[7.5]
Soit $K_v$ un corps local non archim\'edien. On d\'efinit un faisceau de Weil sur $\Spec(K_v)$ en paraphrasant 7.3. Soit $\overline{K}_v$ une cl\^oture s\'eparable de $K_v$. Il revient au m\^eme de se donner un faisceau (de Weil) de $\Lambda$-modules constructible $\mathcal{G}$ sur $\Spec(K_v)$ ou de se donner (a) la fibre $\mathcal{G}_{\overline{K}_v}$, un $\Lambda$-$W(\overline{K}_v/K_v)$-module, de type fini sur $\Lambda$; (b) la fibre $\mathcal{G}_{\overline{k(v)}}$, un $\Lambda$-$W(\overline{k(v)}/k(v))$-module, de type fini sur $\Lambda$; (c) la fl\`eche de sp\'ecialisation, $W(\overline{K}_v/K_v)$-\'equivariante $\mathrm{sp}: \mathcal{G}_{\overline{K}_v} \to \mathcal{G}_{\overline{k(v)}}$.

Soient $v$ une place de $K$ et $x_v: \Spec(K_v) \to X$. Par restriction, tout faisceau (de Weil) de $\Lambda$-modules constructible sur $X$ en d\'efinit un sur $x_v$.
\end{definition}

\begin{definition}[7.6]
Un faisceau est plat si ses fibres sont des $\Lambda$-modules plats.
\end{definition}

\begin{definition}[7.6]
Soit $K_v$ un corps local non archim\'edien. On pose $\deg(v) = [k(v):\mathbb{F}_p]$ et on note $\omega_t$ le quasi-caract\`ere non ramifi\'e de $W(\overline{K}_v/K_v)$, \`a valeurs dans $\Lambda(t)^*$, tel que $\omega_t(F_v) = t^{\deg(v)}$.

Soit $\mathcal{G}$ un faisceau constructible plat sur $\Spec(\mathcal{O}_v)$. On pose
\begin{equation}
Z(\mathcal{G}_v,t) = \det(1 - F_v t^{\deg(v)}, \mathcal{G}_{\overline{k(v)}})^{-1} = \det(1 - F_v \cdot t, \mathcal{G}_{\overline{k(v)}} \otimes \omega_t)^{-1}.
\tag{7.6.0}
\end{equation}
Pour $p$ inversible dans $\Lambda$, on d\'efinit un conducteur
\begin{equation}
a(\mathcal{G}_v) = \sw(\mathcal{G}_{\overline{K}_v}) + \dim(\mathcal{G}_{\overline{K}_v}) - \dim(\mathcal{G}_{\overline{k(v)}})
\tag{7.6.1}
\end{equation}
(un entier si $\Lambda$ est local, une fonction localement constante sur $\Spec(\Lambda)$ en g\'en\'eral).

Supposons que $\Lambda$ soit un corps de caract\'eristique $\not= p$. Soient $\psi$ un caract\`ere additif non trivial de $K_v$ \`a valeurs dans $\Lambda^*$, et $dx$ une mesure de Haar \`a valeurs dans $\Lambda$, comme en 6.5. On pose
\begin{equation}
\varepsilon(\mathcal{G}_v,t) = \varepsilon_0(\mathcal{G}_{\overline{K}_v} \otimes \omega_t, \psi, dx) \cdot \det(-F_v t^{\deg(v)}, \mathcal{G}_{\overline{k(v)}})^{-1}.
\tag{7.6.2}
\end{equation}
$\varepsilon$ est un mon\^ome. D'apr\`es (5.5.2), son degr\'e en $t$ est $a(\mathcal{G}_v) \cdot \deg(v)$.
\end{definition}

\begin{remark}[7.7]
Soit $M$ une repr\'esentation de $W(\overline{K}_v/K_v)$. En termes galoisiens, on d\'efinit des faisceaux $j_!M$ et $j_*M$ sur $\Spec(\mathcal{O}_v)$ par les fl\`eches de sp\'ecialisation $0 \to M$ et $M \xrightarrow{\Id} M$. Les quantit\'es $\varepsilon$ et $\varepsilon_0$ des \S\ pr\'ec\'edents s'identifient \`a la quantit\'e (7.6.2) relative respectivement \`a $j_*M$ et $j_!M$. Au contraire de $j_*$, le foncteur $j_!$ commute \`a la r\'eduction mod $\ell$. Ceci peut expliquer pourquoi, au \S\ 6, nous n'avons d\'efini que $\varepsilon_0$.
\end{remark}

\begin{definition}[7.8]
Soit $\mathcal{G}$ un faisceau constructible de $\Lambda$-modules plats sur $X$. Pour chaque place $v$ de $K$, $\mathcal{G}$ induit un faisceau $\mathcal{G}_v$ sur $\Spec(\mathcal{O}_v)$. On pose
\begin{equation}
Z(\mathcal{G},t) = \prod_v Z(\mathcal{G}_v,t) \in \Lambda[[t]].
\tag{7.8.0}
\end{equation}
\end{definition}

\begin{definition}[7.9]
Supposons que $\Lambda$ soit un corps de caract\'eristique $\ell \not= p$. Il existe sur l'anneau des ad\`eles de $K$ une unique mesure de Haar \`a valeurs dans $\mathbb{Z}[1/p]$ telle que $\int_{\mathbb{A}/K} dx = 1$. On peut \'ecrire $dx = \otimes dx_v$ ($dx_v =$ mesure de Haar sur $K_v$, \`a valeurs dans $\mathbb{Z}[1/p]$; pour presque tout $v$, $\int_{\mathcal{O}_v} dx_v = 1$). Les $dx_v$ d\'efinissent des mesures de Haar \`a valeurs dans $\Lambda$.

S'il existe un caract\`ere non trivial $\psi$ de $\mathbb{A}/K$, \`a valeurs dans $\Lambda^*$, on pose
\begin{equation}
\varepsilon(\mathcal{G},t) = \prod_v \varepsilon(\mathcal{G}_v,t).
\tag{7.9.0}
\end{equation}
Presque tous les termes du produit sont \'egaux \`a 1, et le produit ne d\'epend pas du choix de $\psi$ ni de la d\'ecomposition choisie de $dx$.

Soit $\Lambda'$ d\'eduit de $\Lambda$ par adjonction d'une racine primitive $p$-i\`eme de l'unit\'e. Il existe $\psi$ \`a valeurs dans $\Lambda'^*$. La constante globale correspondante, \'etant ind\'ependante de $\psi$, est invariante par $\Gal(\Lambda'/\Lambda)$, donc dans $\Lambda$. Ceci permet de d\'efinir $\varepsilon$ m\^eme s'il n'existe pas de caract\`ere $\psi$.
\end{definition}

\begin{definition}[7.10]
Notons $i$ l'inclusion de $\Spec(K)$ dans $X$. Si $V$ est une repr\'esentation de $W(\overline{K}/K)$, on note $i_*V$ le faisceau sur $X$ de fibre en $\overline{K}$ \'egale \`a $V$, et pour lequel les $\mathrm{sp}_v$ sont $\mathrm{sp}_v: V^{I_v} \hookrightarrow V$.
\end{definition}


\begin{theorem}[7.11]
Soit $\Lambda$ un corps de caract\'eristique $\ell \not= p$. Soient $\mathcal{G}$ un faisceau constructible de $\Lambda$-vectoriels, $V$ la fibre de $\mathcal{G}$ en $\overline{K}$.
\begin{enumerate}
\item[(i)] $Z(\mathcal{G},t)$ est une fraction rationnelle en $t$.
\item[(ii)] Pour $t \to \infty$, $Z(\mathcal{G},t)$ est asymptotique \`a $\varepsilon(\mathcal{G},t)$. En d'autres termes, $Z(\mathcal{G},t)/\varepsilon(\mathcal{G},t)$ est une s\'erie formelle en $t^{-1}$, de terme constant 1.
\item[(iii)] $Z(\mathcal{G},t) = \varepsilon(\mathcal{G},t) \cdot Z(D(\mathcal{G}),t^{-1})$.
\end{enumerate}
\end{theorem}

\begin{proof}
Soit $S$ un ensemble fini de places contenant toutes les places o\`u la repr\'esentation $V$ est ramifi\'ee. Soit $j$ l'inclusion de $X - S$ dans $X$, et soit $j_!V$ le faisceau sur $X$, de fibre g\'en\'erale $V$, non ramifi\'e en dehors de $S$, et de fibre nulle en $v \in S$.

On v\'erifie aussit\^ot que (i) (ii) \'equivalent \`a
\begin{enumerate}
\item[(i')] $Z(j_!V,t)$ est fonction rationnelle de $t$.
\item[(ii')] Pour $t \to \infty$, $Z(j_!V,t)$ est asymptotique \`a $\varepsilon(j_!V,t)$.
\end{enumerate}

Supposons que $\ell \not= 0$. Pour toute repr\'esentation $W$ de $W(\overline{K}/K)$, non ramifi\'e en dehors de $S$ (identifi\'ee \`a un faisceau localement constant sur $X - S$), posons $j_!^g W = i_*W$, et $R^1j_!^g W =$ le faisceau nul en dehors de $S$, de fibre en $v \in S$ \'egale \`a $H^1(I_v,W)$.

\begin{equation}
Z(Rj_!^g W,t) = \prod_{v \notin S} \det(1-F_v t^{\deg(v)}, W_v)^{-1} \cdot \prod_{v \in S} \frac{\det(1-F_v t^{\deg(v)}, H^0(I_v,W))}{\det(1-F_v t^{\deg(v)}, H^1(I_v,W))}
\tag{7.11.0}
\end{equation}
et
\[
\varepsilon(j_!^g W,t) = \prod_{v \notin S} \varepsilon(W_v,\psi_v,dx_v,t) \cdot \prod_{v \in S} \varepsilon_0(W_v,\psi_v,dx_v,t).
\]

Soit $P' = \Ker(t_\ell: I_v \to \mathbb{Z}_\ell(1))$. Pour toute repr\'esentation $W$ de $I_v$ sur $\Lambda$, on a canoniquement
\begin{equation}
H^0(I_v,W) = W^{P'} \cap W^{\mathbb{Z}_\ell(1)} \quad (= W^{I_v})
\tag{7.11.1}
\end{equation}
Pour toute repr\'esentation $W$ de $\mathbb{Z}_\ell(1)$, on a canoniquement
\begin{align}
H^0(\mathbb{Z}_\ell(1),W) &= W^{\mathbb{Z}_\ell(1)} \quad \text{(invariants)} \tag{7.11.2}\\
H^1(\mathbb{Z}_\ell(1),W) &= W_{\mathbb{Z}_\ell(1)}(-1) \quad \text{(coinvariants tordu).} \notag\\
H^i(\mathbb{Z}_\ell(1),W) &= 0 \quad (i \geq 2). \notag
\end{align}
Puisque la dualit\'e \'echange invariants et coinvariants, on a (cf. 2.3):

\begin{lemma}[7.11.4]
Pour $\ell \not= 0$, $H^0(I_v,W)$ et $H^1(I_v,W)$ sont canoniquement en dualit\'e.
\end{lemma}
Nous n'utiliserons pas que cette dualit\'e peut s'interpr\'eter comme un cup-produit \`a valeurs dans $\mu_\ell$, mod $\ell$.

\begin{equation}
H^1(I_v,\Lambda(1)) = \Hom(I_v,\Lambda(1)) \cong \Lambda,
\tag{7.11.3}
\end{equation}
engendr\'e par un g\'en\'erateur $a$.

Pour $\ell \not= 0$, dans le courant de cette d\'emonstration, nous poserons et d\'efinirons $Z(Rj_!^g W,t)$ par (7.11.0). Cette d\'efinition est justifi\'ee par 7.11.5 et 7.11.6 ci-dessous.

Soient $\Lambda$ un anneau de valuation discr\`ete d'in\'egale caract\'eristique $\ell \not= p$ et $W$ un $\Lambda$-module libre sur lequel $W(\overline{K}_v/K_v)$ agit. Soient $\eta$ le corps des fractions de $\Lambda$, $s = \Lambda/\mathfrak{m}$ le corps r\'esiduel.

\begin{lemma}[7.11.5]
$Z_v(Rj_!^g W_\eta,t) \equiv Z_v(Rj_!^g W_s,t) \pmod{\mathfrak{m}}$.
\end{lemma}

La d\'emonstration ci-dessous perdra son myst\`ere au \S\ 10.

Soient $P' = \Ker(t_\ell: I \to \mathbb{Z}_\ell(1))$, $a$ un g\'en\'erateur de $\mathbb{Z}_\ell(1)$, $F$ un Frobenius g\'eom\'etrique et $K$ le complexe concentr\'e en degr\'es 0 et 1
\[
K: \quad W^{P'} \xrightarrow{\;a-1\;} W^{P'}.
\]
L'endomorphisme $\sum_{i=1}^{q-1} a^i$ modulo l'id\'eal de $W^{P'}$ engendr\'e par $(a-1)W^{P'}$ est inversible: il suffit de le voir apr\`es r\'eduction modulo l'id\'eal maximal de $\Lambda$. Apr\`es r\'eduction, puisqu'un $a$ est l'identit\'e, $a$ est unipotent et $q$ est l'unique valeur propre de $\sum a^i$.

Soit $F = (F_0,F_1)$ l'endomorphisme de $K$ de composantes $F$ et $\sum_{i=1}^{q-1} a^i \circ F$. La quantit\'e
\[
Z = \det(1 - F t^{\deg(v)}, W^{P'})^{-1} \cdot \det(1 - F_1 t^{\deg(v)}, W^{P'})
\]
est de formation compatible \`a l'extension des scalaires \`a $\eta$ ou $s$. Apr\`es extension des scalaires \`a $\eta$ (resp. $s$), on a $Z = \det(1 - F t^{\deg(v)}, H^0(K))^{-1} \cdot \det(1 - F t^{\deg(v)}, H^1(K))$. On conclut en notant que $H^0(K) = H^0(I_v,W)$ et $H^1(K) = H^1(I_v,W)$.

\textbf{Variante 7.11.6.} Supposons $\Lambda$ complet et soit $W$ un $\Lambda$-module sur lequel $W(\overline{K}_v/K_v)$ agit contin\^ument pour la topologie $\ell$-adique (et non plus, comme auparavant, pour la topologie discr\`ete). On pose encore $H^1(I_v,W) = W_{P'}^{\mathbb{Z}_\ell(1)}(-1)$ (coinvariants tordu), et on d\'efinit $Z_v(Rj_!^g W,t)$ comme auparavant. Le $H^1$ d\'efini plus haut peut s'interpr\'eter comme un $H^1$ continu
\begin{equation}
H^1(I_v,W) = \varprojlim_n H^1(I_v,W/\ell^n W).
\tag{7.11.6.1}
\end{equation}
La formule de r\'eduction 7.11.5 reste valable (avec la m\^eme d\'emonstration).

\begin{lemma}[7.11.7]
La formule (iii) \'equivaut \`a
\begin{equation}
Z(j_!V,t) = \varepsilon(j_!V,t) \cdot Z(Rj_!^g V^*(1),t^{-1}).
\tag{7.11.7}
\end{equation}
\end{lemma}

7.11.7 r\'esulte de l'identit\'e locale
\[
\det(1 - F_v t^{\deg(v)}, V^{I_v})^{-1} = \varepsilon_{0,v}(V,t) \cdot \det(-F_v t^{\deg(v)}, H^1(I_v,V))^{-1} \cdot \det(1 - F_v t^{-\deg(v)}, H^1(I_v,V^*(1))).
\]
Si on remplace $V$ par la repr\'esentation $V \otimes \omega_t$ de $W(\overline{K}_v/K_v)$ sur $\Lambda((t))$, cette formule se simplifie en
\[
\det(1-F_v,V^{I_v})^{-1} = \varepsilon_{0,v}(V,\psi_v,dx_v) \cdot \det(-F_v,H^1(I_v,V))^{-1} \cdot \det(1-F_v,H^1(I_v,V^*(1)))
\]
qui r\'esulte de la dualit\'e locale 7.11.4.

Prouvons 7.11. Si $\Lambda = \mathbb{C}$, le produit infini $Z(j_!V,t)$ converge pour $|t|$ assez petit, et $L(V,s)$ est une fonction m\'eromorphe de $t$; l'\'equation fonctionnelle montre qu'elle est \`a croissance mod\'er\'ee pour $t \to \infty$, donc une fonction rationnelle de $t$. L'\'equation fonctionnelle de la fonction $L$ fournit (iii). Les \'enonc\'es (i) et (iii) sont purement alg\'ebriques. \'Etant vrais pour $\Lambda = \mathbb{C}$ (ainsi que (i') et (7.11.7)), ils restent vrais pour $\Lambda$ de caract\'eristique 0.

Supposons maintenant $\Lambda$ de caract\'eristique $\ell \not= 0$. Soit $J$ un sous-groupe ouvert distingu\'e de $W(\overline{K}/K)$, contenant $I_v$ pour $v \in S$. Si une repr\'esentation $V$ de $W(\overline{K}/K)/J$ sur $\Lambda$ se rel\`eve en caract\'eristique 0, les assertions (i') et (7.11.7) pour $V$ s'obtiennent par r\'eduction mod $\ell$ (compte tenu, pour (7.11.7), de 7.11.5). Les quantit\'es $Z(j_!V,t)$, $Z(Rj_!^g V,t)$ et $\varepsilon(j_!V,t)$ d\'ependent additivement de $V$ (une extension donne un produit; cela r\'esulte de (7.11.6.1)). Pour obtenir (i) et (7.11.7) en g\'en\'eral il suffit donc de savoir que, apr\`es extension des scalaires toute repr\'esentation de $W(\overline{K}/K)/J$ se rel\`eve virtuellement en caract\'eristique 0. La m\'ethode de 4.10 (cf. 4.10.3) ram\`ene cette question au cas des groupes finis, i.e. au th\'eor\`eme de Brauer.

L'\'enonc\'e (ii') r\'esulte de (7.11.7) et ce que, pour $t \to 0$, $Z \to 1$.
\end{proof}

\begin{proposition}[7.12]
Si $\Lambda$ est un corps de caract\'eristique $p$, $\mathcal{G}$ un faisceau constructible de $\Lambda$-vectoriels, $Z(\mathcal{G},t)$ est encore une fraction rationnelle en $t$.
\end{proposition}
Dans la d\'emonstration de 7.11 (i) nous n'avons en effet pas fait usage de l'hypoth\`ese $\ell \not= p$.


\section{Repr\'esentations $\ell$-adiques des groupes de Weil locaux non archim\'ediens. Structures de Hodge et groupes de Weil locaux archim\'ediens}

\begin{definition}[8.1]
Soit $K$ un corps local non archim\'edien. Nous utiliserons les notations 2.1 et 2.2. Soient $\ell$ un nombre premier $\not= p$, $E_\lambda$ une extension finie de $\mathbb{Q}_\ell$ (on pourrait prendre pour $E_\lambda$, plus g\'en\'eralement, un corps valu\'e extension de $\mathbb{Q}_\ell$). Une repr\'esentation $\lambda$-adique $V$ de $W(\overline{K}/K)$ est une repr\'esentation continue pour la topologie $\lambda$-adique $\rho: W(\overline{K}/K) \to \GL(V)$ de $W(\overline{K}/K)$ sur un $E_\lambda$-vectoriel de dimension finie.
\end{definition}

Le r\'esultat suivant est d\'emontr\'e dans [13].

\begin{theorem}[8.2 --- Grothendieck]
Soit $(V,\rho)$ une repr\'esentation $\lambda$-adique de $W(\overline{K}/K)$. Il existe $N \in \End(V)(-1)$, nilpotent, tel que pour $\sigma$ dans un sous-groupe d'indice fini de $I$, on ait $\rho(\sigma) = \exp(N \cdot t_\ell(\sigma))$. Ce $N \in \End(V)(-1)$ est unique, donc invariant par Galois: on a pour $w \in W(\overline{K}/K)$, $\rho(w) N \rho(w)^{-1} = q^{v'(w)} \cdot N$.
\end{theorem}

Nous allons utiliser 8.2 pour d\'ecrire les repr\'esentations $\lambda$-adiques de $W(\overline{K}/K)$ en termes ind\'ependants de la topologie de $E_\lambda$. Ce sera fait deux fois: 8.3 est plus intrins\`eque, et 8.4, qui nous suffit, plus \'el\'ementaire.

\begin{definition}[8.3.1]
\`A isomorphisme unique pr\`es, il existe un et un seul groupe alg\'ebrique $\mathbb{G}$ sur $\mathbb{Q}_\ell$, isomorphe \`a $\mathbb{G}_a$, et muni d'un isomorphisme $\mathbb{Q}_\ell(1) \cong \mathbb{G}(\mathbb{Q}_\ell)$. On le note $\mathbb{G}_a(1)$.
\end{definition}

\begin{definition}[8.3.2]
Soit $J \subset I$ un sous-groupe invariant ouvert de $W(\overline{K}/K)$, et $J' = J \cap \Ker(t_\ell)$. Le quotient $W(\overline{K}/K)/J'$ est extension du groupe discret $W(\overline{K}/K)/J$ par $J/J'$; appliquant \`a cette extension $t_\ell: J/J' \hookrightarrow \mathbb{Z}_\ell \cong \mathbb{G}_a(1)$, on obtient une extension $W^J$ de $W(\overline{K}/K)/J$ par $\mathbb{G}_a(1)$ (voir Serre [12] II 1.3).
\[
0 \to \mathbb{G}_a(1) \to W^J \to W(\overline{K}/K)/J \to 0
\]
\end{definition}

\begin{definition}[8.3.3]
Les $W^J$ forment un syst\`eme projectif, et 8.2 signifie que toute repr\'esentation $\lambda$-adique de $W(\overline{K}/K)$ se factorise par une repr\'esentation alg\'ebrique d'un $W^J$, i.e. par une repr\'esentation du sch\'ema en groupes $\varprojlim W^J$.
\end{definition}

\begin{definition}[8.3.4]
Soient $\widetilde{I}$ l'image r\'eciproque de $I$ dans $\widetilde{W}(\overline{K}/K)$ et $\widetilde{I}'' = \varprojlim_J \widetilde{I}/[\widetilde{I},J]$. On a un diagramme commutatif
\[
\begin{CD}
0 @>>> I @>>> W(\overline{K}/K) @>>> \mathbb{Z} @>>> 0 \\
@. @VVV @VVV @VVV @. \\
0 @>>> I \times \mathbb{G}_a(1) @>>> W_L(\overline{K}/K) @>>> \mathbb{Z} @>>> 0 \\
@. @VVV @VVV @VVV @. \\
0 @>>> \mathbb{G}_a(1) @>>> \mathbb{G}_a(1) @>>> 0 @>>> 0
\end{CD}
\]
\end{definition}

Soit $'W_\ell$ le produit semi-direct de $W(\overline{K}/K)$ par $\mathbb{G}_a(1)$, sur lequel $W(\overline{K}/K)$ agit par $w x w^{-1} = q^{v'(w)} \cdot x$. Il existe des isomorphismes entre $W(\overline{K}/K)$ et $'W_\ell(\overline{K}/K)$ qui rendent commutatif
\[
\begin{CD}
I @>>> W(\overline{K}/K) @>>> \mathbb{Z} \\
@VVV @VVV @VV{=}V \\
I \times \mathbb{G}_a(1) @>>> 'W_\ell(\overline{K}/K) @>>> \mathbb{Z}
\end{CD}
\]
et on v\'erifie facilement que deux quelconques sont conjugu\'es par un \'el\'ement de $\mathbb{G}_a(1)$ (l'ind\'etermination est dans le choix d'un rel\`evement d'un Frobenius $F \in W(\overline{K}/K)$, et on utilise que $q^{-1} \not= 0$; voir 8.4.3 ci-dessous).

\begin{definition}[8.3.6]
Soit $'\mathfrak{W}$ le sch\'ema en groupes sur $\mathbb{Q}$, produit semi-direct de $W(\overline{K}/K)$ par $\mathbb{G}_a$, sur lequel $W(\overline{K}/K)$ agit par $w x w^{-1} = q^{-v'(w)} \cdot x$. Chaque isomorphisme $\mathbb{G}_a \cong \mathbb{G}_a(1)$ (i.e. $\mathbb{Q} \cong \mathbb{Q}(1)$) d\'efinit un isomorphisme $\ '\mathfrak{W} \cong 'W_\ell$, d'o\`u une classe de tels automorphismes modulo $\mathbb{Q}^*$ agissant sur $'W_\ell$ via son action sur $\mathbb{G}_a(1)$.

On v\'erifie (voir 8.4.3 ci-dessous) que ce groupe d'automorphismes agit trivialement sur l'ensemble des classes d'isomorphismes de repr\'esentations de $'W_\ell$. Au total, on a obtenu que
\begin{equation}
\text{Les classes d'isomorphie de repr\'esentations $\lambda$-adiques de $W(\overline{K}/K)$ sont en correspondance bijective naturelle avec}
\tag{8.3.7}
\end{equation}
les classes d'isomorphie de repr\'esentations sur $E_\lambda$ du sch\'ema en groupes sur $\mathbb{Q}$ $\ '\mathfrak{W}$.
\end{definition}

\begin{definition}[8.4.1]
Soit $E$ un corps de caract\'eristique 0. Une repr\'esentation de $\ '\mathfrak{W}(\overline{K}/K)$ sur $E$ est un couple $\sigma = (\rho',N')$ consistant en
\begin{enumerate}
\item[(a)] une repr\'esentation de dimension finie $\rho': W(\overline{K}/K) \to \GL(V)$ de $W(\overline{K}/K)$ sur $E$;
\item[(b)] un endomorphisme nilpotent $N'$ de $V$, tel que
\begin{equation}
\rho'(w) N' \rho'(w)^{-1} = q^{v'(w)} \cdot N'.
\tag{8.4.1.1}
\end{equation}
\end{enumerate}
\end{definition}

\begin{definition}[8.4.2]
Choisissons un Frobenius g\'eom\'etrique $F \in W(\overline{K}/K)$ et un isomorphisme $\tau: \mathbb{Q}_\ell(1) \cong \mathbb{Q}_\ell$. On associe comme suit une repr\'esentation $(\rho',N')$ de $\ '\mathfrak{W}(\overline{K}/K)$ sur $E_\lambda$ \`a une repr\'esentation $\lambda$-adique $(V,\rho)$ de $W(\overline{K}/K)$:
\begin{enumerate}
\item[(a)] Pour $\sigma \in I$, on pose $\rho'(\sigma) = \rho(\sigma) \cdot \exp(-N \cdot \tau^{-1} \circ t_\ell(\sigma))$ ($N$ \'etant d\'efini par 8.2).
\item[(b)] $N' \in \End(V)$ correspond \`a $N$, via $\tau$.
\end{enumerate}
\end{definition}

\begin{lemma}[8.4.3]
La classe d'isomorphie de $(\rho',N')$ ne d\'epend que de $\rho$, non des choix de $F$ et $\tau$.
\end{lemma}

\begin{proof}
Changeons $F$ en $F\beta$, pour $\beta \in I$, pour obtenir $(\rho'',N')$. On a pour $\lambda = \exp((1-q^{-1})^{-1} t_\ell(\beta) \cdot N)$: $\rho'' = \lambda \rho' \lambda^{-1}$, $N' = \lambda N' \lambda^{-1}$.

Changer $\tau$ en $\delta\tau$ change $(\rho',N')$ en $(\rho',\delta N')$. Nous montrerons que pour toute repr\'esentation $(\rho',N')$ de $\ '\mathfrak{W}(\overline{K}/K)$ sur un corps $E$, et tout $\alpha \in E$, $(\rho',N') \cong (\rho',\alpha N')$.

Soient $I' = \Ker(\rho')$ et $n$ tel que $F^n$ soient central dans $W/I'$ (par exemple, $n \geq |\Aut(I/I')|$). Pour $\alpha$ une classe de conjugaison dans une cl\^oture alg\'ebrique de $E$, soit $V_\alpha$ le plus grand sous-espace de $V$, stable par $F^n$, tel que les valeurs propres de $F^n|V_\alpha$ soient dans $\alpha$. On v\'erifie que
\begin{enumerate}
\item[(a)] $V_\alpha$ est stable sous $\rho'(W(\overline{K}/K))$, et $V = \bigoplus V_\alpha$;
\item[(b)] $N': V_\alpha \to V_{q^{-n}\alpha}$.
\end{enumerate}
Si $A$ est la transformation lin\'eaire, laissant stable les $V_\alpha$, et telle que $A|V_\alpha = A(\alpha)$, avec $A(q^{-n}\alpha) = \alpha \cdot A(\alpha)$, on a donc $A(\rho',N')A^{-1} = (\rho',\alpha N')$.
\end{proof}

De 8.4.3 on d\'eduit aussit\^ot la conclusion 8.3.7.

\begin{definition}[8.5]
Soit $(\rho',N')$ une repr\'esentation de $\ '\mathfrak{W}(\overline{K}/K)$ sur $E$. Soit $u$ la composante unipotente de $\rho'(F)$: $\rho'(F) = \rho'(F)^{\mathrm{ss}} \cdot u$. Puisque $N' \in \End(V)$ est vecteur propre de $\rho'(F)$, $u$ et $N'$ commutent. Puisqu'une puissance de $F$ est centrale dans $\mathrm{Im}(\rho')$, une puissance de $u$, donc $u$, centralise $\mathrm{Im}(\rho')$. Si $w_1, w_2 \in W(\overline{K}/K)$, $\rho'(w_1)$ et $\rho'(w_2)$ ont une puissance commune. On en d\'eduit que $u$ est la composante unipotente de $\rho'(F^n\sigma)$ ($\sigma \in I$), et, pour $n \not= 0$, aussi celle de $\rho'(F^n\sigma) \cdot \exp((1-q^{-n})^{-1} \cdot a \cdot N')$, qui en est le conjugu\'e par $\exp((1-q^{-n})^{-1} \cdot a \cdot N')$. Posons
\begin{equation}
\nu = u \cdot \exp(-(1-q^{-n})^{-1} \cdot a \cdot N').
\tag{8.5.1}
\end{equation}
\end{definition}

\begin{definition}[8.6]
\begin{enumerate}
\item[(i)] $(\rho'^{\mathrm{ss}},N')$ est la $F$-semi-simplifi\'ee de $(\rho',N')$.
\item[(ii)] On dit que $(\rho',N')$ est $F$-semi-simple si $\rho' = \rho'^{\mathrm{ss}}$, i.e. si pour un (resp. pour tout) $w \in W(\overline{K}/K) - I$ (resp. $\in I$), $\rho'(w) \cdot \exp(aN')$ est semi-simple.
\end{enumerate}
Pour une repr\'esentation $\lambda$-adique, on d\'efinit encore sa $F$-semi-simplifi\'ee par (8.5.1), et la $F$-semi-simplicit\'e comme en (ii).
\end{definition}

\begin{definition}[8.7]
Le formalisme 8.3 ou 8.4 permet de comparer des repr\'esentations $\lambda$-adiques pour diff\'erents corps $E_\lambda$. Pour ce faire, introduisons la terminologie classique suivante:

Soit $F$ une extension de $E$, et $\sigma$ une repr\'esentation de $\ '\mathfrak{W}(\overline{K}/K)$ sur $F$. On dit que $\sigma$ est d\'efinie sur $E$ si, pour une (pour une) extension alg\'ebriquement close $\overline{F}$ de $F$, $\sigma_{\overline{F}}$ est isomorphe \`a ses conjugu\'es sous $\Aut(\overline{F}/E)$.

Soient $F_1$ et $F_2$ deux extensions de $E$, et $\sigma_i$ une repr\'esentation de $\ '\mathfrak{W}(\overline{K}/K)$ sur $F_i$. On dit que $\sigma_1$ et $\sigma_2$ sont compatibles si elles sont d\'efinies sur $E$ et que, pour $\overline{F} \supset F_1, F_2 \supset E$ une extension alg\'ebriquement close commune, $\sigma_{1\overline{F}}$ et $\sigma_{2\overline{F}}$ sont isomorphes.
\end{definition}

\begin{definition}[8.8]
Soit $E$ un corps de nombres, $\lambda_1, \lambda_2$ deux places finies ne divisant pas $p$ de $E$ et $\rho_1, \rho_2$ des repr\'esentations $\lambda_1$- et $\lambda_2$-adiques de $W(\overline{K}/K)$. On dit que $\rho_1$ et $\rho_2$ sont compatibles si les repr\'esentations correspondantes de $\ '\mathfrak{W}(\overline{K}/K)$ le sont (8.7).
\end{definition}

On laisse au lecteur la v\'erification du plaisant exercice suivant.

\begin{proposition}[8.9]
Supposons $(V_1,\rho_1)$ et $(V_2,\rho_2)$ $F$-semi-simples. Soit $W_\bullet$ l'unique filtration finie croissante de $V_i$ telle que $N W_k \subset W_{k-2}$ et que $N$ induise un isomorphisme $\gr_k^W(V_i) \xrightarrow{\;\sim\;} \gr_{-k}^W(V_i)$. Pour que $\rho_1$ et $\rho_2$ soient compatibles, il faut et il suffit que les caract\`eres des repr\'esentations de $W(\overline{K}/K)$ sur les $\gr_k^W(V_1)$ et $\gr_k^W(V_2)$ soient \`a valeurs dans $E$ et respectivement \'egaux.
\end{proposition}

\begin{example}[8.10 --- on fait $E = \mathbb{Q}$]
Soit $A$ une vari\'et\'e ab\'elienne sur $K$. Les repr\'esentations de $W(\overline{K}/K) \subset \Gal(\overline{K}/K)$ sur les $V_\ell(A) = T_\ell(A) \otimes \mathbb{Q}_\ell$ sont $F$-semi-simples et compatibles.
\end{example}

\textbf{Variante 8.11.} Soit $G$ un groupe alg\'ebrique lin\'eaire sur $E$ de caract\'eristique 0. Une repr\'esentation de $\ '\mathfrak{W}$ dans $G$ (sur $E$) est un couple $(\rho',N')$: $\rho': W(\overline{K}/K) \to G(E)$ (trivial sur un sous-groupe ouvert) et $N' \in \Lie(G)$, (8.4.1.1) \'etant v\'erifi\'e. Pour $E = E_\lambda$, les formules (8.4.2) associent une repr\'esentation de $\ '\mathfrak{W}$ dans $G$ \`a toute repr\'esentation continue pour la topologie $\lambda$-adique $\rho: W \to G(E_\lambda)$. Sa classe d'isomorphie g\'eom\'etrique (apr\`es extension des scalaires \`a $\overline{E}_\lambda$) ne d\'epend pas des choix arbitraires faits: la premi\`ere partie de la d\'emonstration 8.4.3 s'applique encore, il reste \`a montrer que $(\rho',N') \sim (\rho',\alpha N')$. Pour ce, on note que le sous-groupe alg\'ebrique $H$ de $G$ qui centralise $\mathrm{Im}(\rho')$ et envoie $N'$ sur un multiple contient un $\rho'(F^m)$, donc un \'el\'ement $g$ tel que $\mathrm{ad}_g(N') = q^m N'$. Le caract\`ere de $H$ donnant l'action sur $N'$ ne peut donc qu'\^etre surjectif.

Si $G$ est d\'efini sur un corps de nombres $E$, on d\'efinit la compatibilit\'e de repr\'esentations $\lambda$-adiques $W(\overline{K}/K) \to G(E_\lambda)$ comme en 8.7, 8.8. Le sorite 8.5, 8.6 se g\'en\'eralise tel quel.

\begin{definition}[8.12 --- Facteurs locaux]
Soit $(\rho',N')$ une repr\'esentation de $\ '\mathfrak{W}(\overline{K}/K)$ sur un $E$-vectoriel $V$. Par abus de notation, on pose $V^I = \Ker(N')^{\rho'(I)}$; c'est une repr\'esentation de $W(\overline{k}/k)$.

On d\'efinit un conducteur, un facteur $L$ local et des constantes locales par les formules
\begin{align}
a(V) &= \sw(V,\rho') + \dim V - \dim V^I \tag{8.12.1}\\
L(V,t) &= \det(1-Ft,V^I)^{-1} \tag{8.12.2}\\
\varepsilon_0(V,t) &= \varepsilon_0(V \otimes \omega_t) \tag{8.12.3}\\
\varepsilon(V,t) &= \varepsilon_0(V,t) \cdot \det(-Ft,V^I)^{-1} \tag{8.12.4}
\end{align}
Pour $(\rho',N')$ d\'efinis par une repr\'esentation $\lambda$-adique $\rho$ sur $V$, on a $V^I = V^{\rho(I)}$, $Z(V,t) = \det(1-Ft,V^{\rho(I)})$, $\varepsilon_0(V,t) = \varepsilon_0(\text{semi-simplifi\'ee de }V,t)$.
\end{definition}

\textbf{Le cas archim\'edien:}

Dans le cas non archim\'edien, ce que fournit la g\'eom\'etrie (la cohomologie $\ell$-adique) sont des repr\'esentations $\lambda$-adiques, i.e. des repr\'esentations de $\ '\mathfrak{W}(\overline{K}/K)$, non des repr\'esentations de $W(\overline{K}/K)$.

Dans le cas archim\'edien, de m\^eme, la g\'eom\'etrie fournit non des repr\'esentations de $W(\overline{K}/K)$, mais des structures de Hodge. Voici la r\`egle \`a utiliser pour passer des structures de Hodge aux repr\'esentations.

\textbf{(A) $K$ complexe.}

La g\'eom\'etrie fournit des vectoriels $V$ sur $K$, munis d'une bigraduation de Hodge $V = \bigoplus V^{p,q}$ (par exemple, pour $X$ une vari\'et\'e alg\'ebrique sur $K$, $H^n(X,K)$ est muni d'une telle d\'ecomposition). On en d\'eduit une repr\'esentation (semi-simple) $\rho$ de $W(\overline{K}/K) = K^*$ sur $V$ en posant
\[
\rho(z) = \bigoplus_{p,q} z^{-p} \bar{z}^{-q} \quad \text{sur } V^{p,q}.
\]

\textbf{(B) $K$ r\'eel.}

La g\'eom\'etrie fournit des vectoriels $V$ sur $K$, munis d'une bigraduation de Hodge $V = \bigoplus V^{p,q}$ et d'une involution $F_\infty$ telle que $F_\infty(V^{p,q}) = V^{q,p}$. Par exemple, pour $X$ une vari\'et\'e alg\'ebrique sur $K$, $\Gal(\overline{K}/K) = \mathbb{Z}/2$ agit sur $X(\overline{K})$ et $H^n(X(\overline{K}),K)$ est muni d'une telle structure.

On en d\'eduit une repr\'esentation (semi-simple) $\rho$ de $W(\overline{K}/K)$ sur $V$ en posant
\[
\rho(z) = \bigoplus_{p,q} z^{-p} \bar{z}^{-q} \text{ sur } V^{p,q} \text{ pour } z \in K^*, \quad \rho(F) = F_\infty.
\]
Ces formules diff\`erent de celles de [3] 1.12 (d\'efigur\'e en outre par une faute ``d'impression'' p.~21 l.12), parce qu'aux places finies j'utilisais dans loco cit. l'inverse de l'isomorphisme de la th\'eorie du corps de classe local utilis\'e ici.

Aux places finies, $H^2(\mathbb{P}^1, \mathbb{Q}_\ell) \cong \mathbb{Q}_\ell(-1)$ correspond au quasi-caract\`ere $\omega_1$. De m\^eme, $H^2(\mathbb{P}^1(\overline{K}),K)$, structure de Hodge de type $(-1,-1)$, avec $F_\infty = -1$ si $K = \mathbb{R}$, correspond au quasi-caract\`ere $\omega_1$ de $W(\overline{K}/K)$.


\section{Syst\`emes compatibles de repr\'esentations $\ell$-adiques}

Pour simplifier l'expos\'e, nous ne consid\'ererons dans ce \S\ que des repr\'esentations de groupes de Galois (plut\^ot que des repr\'esentations de groupes de Weil).

Comme en 7.1, on note $X$ une courbe projective et lisse sur $\mathbb{F}_p$ irr\'eductible sur $\mathbb{F}_p$, de corps de fonctions $K$.

\begin{definition}[9.1]
Soit $E_\lambda$ un corps local extension finie de $\mathbb{Q}_\ell$, avec $\ell \not= p$. Soit $W$ une repr\'esentation $\lambda$-adique de $\Gal(\overline{K}/K)$, i.e. une repr\'esentation continue pour la topologie de $E_\lambda$ et presque partout non ramifi\'ee de $\Gal(\overline{K}/K)$ dans un vectoriel de dimension finie sur $E_\lambda$. Le groupe de Galois \'etant compact, il existe dans $W$ un r\'eseau invariant $W^0$ (un r\'eseau est un sous-$\mathcal{O}_\Lambda$-module libre de $W$ tel que $E_\lambda \otimes_{\mathcal{O}_\Lambda} W^0 = W$). Les facteurs locaux $\det(1 - F_v t^{\deg(v)}, W_v^{I_v})$ sont donc dans $\mathcal{O}_\Lambda[[t]]$, ainsi que leur produit $Z(W,t)$. Soit $\Gal_0(\overline{K}/K)$ le noyau de l'application de $\Gal(\overline{K}/K)$ dans $\Gal(\overline{\mathbb{F}}_p/\mathbb{F}_p)$. D'apr\`es un th\'eor\`eme de Grothendieck, $Z(W,t)$ est une fraction rationnelle, est m\^eme un polyn\^ome si $W^{\Gal_0(\overline{K}/K)} = 0$ et v\'erifie une \'equation fonctionnelle
\begin{equation}
Z(W,t) = \varepsilon^{\mathrm{Gr}} \cdot Z(W^*, p^{-1}t^{-1})
\tag{9.1.0}
\end{equation}
o\`u $\varepsilon^{\mathrm{Gr}}$ est un mon\^ome de degr\'e $\sum_v \deg(v) \cdot a(i^*W_v)$ (8.12.0).

Le lecteur trouvera au \S\ 10 un r\'esum\'e de la th\'eorie de Grothendieck.
\end{definition}

\begin{definition}[9.2]
Soient $E$ un corps de nombres, $\mathcal{L}$ un ensemble infini de places finies ne divisant pas $p$ de $E$ et $\mathcal{O} \subset E$ l'anneau des \'el\'ements de $E$ $\lambda$-entiers pour tout $\lambda \in \mathcal{L}$. On d\'esigne encore par $\lambda$ l'id\'eal premier de $\mathcal{O}$ d\'efini par $\lambda$ (la valuation $\lambda$-adique).

Supposons donn\'e, pour chaque $\lambda \in \mathcal{L}$, une repr\'esentation $_\lambda V$ de $\Gal(\overline{K}/K)$. On suppose que, pour toute place $v$ de $K$, les repr\'esentations $\lambda$-adiques de $W(\overline{K}_v/K_v)$ soient compatibles, au sens strict (8.8).

Il nous suffira en fait de savoir que leurs $F$-semi-simplifi\'ees (8.6) sont compatibles (cf. aussi 9.8). La fraction rationnelle $Z(_\lambda V,t) \in \mathcal{O}(t)$ et la ``constante'' (un mon\^ome) $\prod_v \varepsilon(_\lambda V_v,\psi_v,dx_v,t) \in \mathcal{O}(t)$ sont donc ind\'ependantes de $\lambda$ (comme d'habitude, les $\psi_v$ sont les composantes d'un caract\`ere non trivial de $\mathbb{A}/K$ et $\otimes dx_v$ est une mesure de Tamagawa).

On note ces fonction et constante $Z(V,t)$ et $\varepsilon(V,t)$.
\end{definition}

\begin{theorem}[9.3]
Sous les hypoth\`eses pr\'ec\'edentes,
\begin{equation}
Z(V,t) = \varepsilon(V,t) \cdot Z(V^*, pt^{-1}).
\tag{9.3.1}
\end{equation}
\end{theorem}

\begin{proof}
Puisque $\mathcal{L}$ est infini, l'application $\mathcal{O} \to \prod_\lambda \mathcal{O}/\lambda$ est injective: il suffit de prouver 9.3.1 apr\`es r\'eduction mod $\lambda$ pour tout $\lambda \in \mathcal{L}$. Nous allons pour cela mettre 9.3.1 sous une forme compatible \`a la r\'eduction mod $\lambda$, et justifiable, mod $\lambda$, de (7.11.8).

Soient $S$ un ensemble fini de places de $K$, contenant toutes les places o\`u $V$ se ramifie et $j$ l'inclusion de $X - S$ dans $X$. Pour toute repr\'esentation $\lambda$-adique $W$, non ramifi\'e en dehors de $S$, nous poserons
\begin{align*}
Z(j_!^g W,t) &= \prod_{v \notin S} \det(1-F_v t^{\deg(v)}, W_v)^{-1} \cdot \prod_{v \in S} \frac{\det(1-F_v t^{\deg(v)}, H^0(I_v,W))}{\det(1-F_v t^{\deg(v)}, H^1(I_v,W))} \\
\varepsilon(j_!^g W,t) &= \prod_{v \notin S} \varepsilon(W_v,\psi_v,dx_v,t) \cdot \prod_{v \in S} \varepsilon_0(W_v,\psi_v,dx_v,t).
\end{align*}

Un calcul facile, d\'ej\`a fait en 7.11.7, montre que (9.3.1) \'equivaut \`a
\begin{equation}
Z(j_!^g(_\lambda V),t) = \varepsilon(j_!^g(_\lambda V),t) \cdot Z(Rj_!^g(_\lambda V^*),t^{-1}).
\tag{9.3.2}
\end{equation}
On conclut par 7.11.6 et 7.11.8.
\end{proof}

\textbf{Stabilit\'e 9.4.} Soient $F$ une extension de $E$ et $\mathcal{L}_F$ l'image r\'eciproque de $\mathcal{L}$. Un syst\`eme compatible de repr\'esentations $\lambda$-adiques $(_\lambda V)_{\lambda \in \mathcal{L}}$ d\'efinit par extension des scalaires un syst\`eme compatible de repr\'esentations $\lambda$-adiques $(_\lambda V)_{\lambda \in \mathcal{L}_F}$. Si $\chi$ est un caract\`ere du groupe des classes d'id\`eles, \`a valeurs dans les racines de l'unit\'e de $F$, les $_\lambda V \cdot \chi$ forment encore un syst\`eme compatible, justiciable de 9.3.

\begin{corollary}[9.5]
Soit $\chi$ un caract\`ere du groupe des classes d'id\`eles, \`a valeurs dans les racines de l'unit\'e d'une extension finie de $E$ (cf. 9.4). Soient $S(\chi)$ (resp. $S(V)$) l'ensemble des places de $K$ o\`u $\chi$ (resp. $V$) se ramifie. Supposons que $S(V) \cap S(\chi) = \emptyset$. On a
\[
\varepsilon((V - \dim V[1]) \otimes (\chi[\chi] - [1]),t) = \prod_{v \in S(\chi)} \varepsilon(V_v,t) \cdot \varepsilon(V,t)^{-1} \cdot \varepsilon(\chi,t)^{-\dim(V)} \cdot \varepsilon([1],t)^{\dim(V)} \cdot a(\chi)^{\dim V} \cdot \det V(\mathfrak{D}_v).
\]
\end{corollary}
C'est un corollaire de 5.5.3.

\begin{example}[9.6 --- pour $E = \mathbb{Q}$, $\mathcal{L} =$ les nombres premiers autres que $p$]
Si $E$ est une courbe elliptique sur $K$, les $H^1(E,\mathbb{Q}_\ell)$ forment un syst\`eme compatible de repr\'esentations $\ell$-adiques. Si l'invariant modulaire n'est pas constant, pour tout caract\`ere $\chi$ comme en 9.5, les fractions rationnelles $Z(E_\chi,t)$ correspondantes sont des polyn\^omes, et v\'erifient 9.5.
\end{example}

\begin{theorem}[9.8]
Soient $\lambda, \mu$ des places de $E$ et $_\lambda V$, $_\mu V$ des repr\'esentations $\lambda$-, $\mu$-adiques de $\Gal(\overline{K}/K)$. Soit $S$ un ensemble fini de places de $K$. Pour $v \notin S$, supposons que les semi-simplifi\'ees de $_\lambda V_v$ et $_\mu V_v$ soient compatibles. Alors, pour toute place $v$, les semi-simplifi\'ees de $_\lambda V_v$ et $_\mu V_v$ sont compatibles.
\end{theorem}

Si $S$ est pris assez grand pour contenir toutes les places ramifi\'ees pour $_\lambda V$ ou $_\mu V$, l'hypoth\`ese se r\'eduit \`a: pour $v \notin S$, les polyn\^omes caract\'eristiques de $F_v$ agissant dans $_\lambda V_v$ et $_\mu V_v$ sont dans $E[t]$ et \'egaux.

Le th\'eor\`eme 9.8 est donc une r\'eponse partielle, dans le cas le moins int\'eressant des corps de fonctions, \`a la question de Serre [12] I 12.

\begin{proof}
On prendra garde aux points suivants.
\begin{enumerate}
\item[(a)] La conclusion porte seulement sur les semi-simplifi\'ees des repr\'esentations locales, et sur leur $F$-semi-simplifi\'ee. En d'autres termes, avec les notations de 8.4.2, on consid\`ere seulement la semi-simplifi\'ee de $\rho'$, et pas $N'$.
\item[(b)] Pour $\lambda, \mu \mid \ell$, $_\lambda V$ et $_\mu V$, le th\'eor\`eme affirme que les semi-simplifi\'ees des repr\'esentations locales $_\lambda V_v$ de $W(\overline{K}_v/K_v)$ sont d\'efinies sur $E$, i.e. ont un caract\`ere \`a valeurs dans $E$.
\end{enumerate}

On peut supposer, et on suppose que $S$ contient toutes les places o\`u $I_v$ agit via un groupe infini. Pour $v \notin S$, il n'y a d\`es lors plus \`a distinguer entre semi-simplifi\'ee de $_\lambda V_v$ et $F$-semi-simplifi\'ee.

Appliquons le th\'eor\`eme de Grothendieck rappel\'e en 9.1. Cette identit\'e se r\'ecrit comme produit de facteurs locaux. Le premier membre est dans $E(t)$ et le m\^eme pour $_\lambda V$ et $_\mu V$. Le second membre est donc dans $E(T)$, ind\'ependant de $\lambda$.

D\'ecomposant le second membre en le produit d'un mon\^ome et d'une fraction rationnelle valant 1 en $t=0$, on trouve par le calcul local de 7.11.7 que le produit local est le m\^eme pour $_\lambda V$ et $_\mu V$. Les facteurs de ce produit sont additifs en $_\lambda V_v$ (une extension donne un produit; cela r\'esulte de (7.11.6.1)), donc les m\^emes pour $_\lambda V_v$ et sa semi-simplifi\'ee. Pour celle-ci, invariants et coinvariants de $I_v$ sont pareils, et on trouve que
\begin{equation}
\prod_{v \in S} \det(1 - F_v t, {_\lambda V_v}^{I_v}) / \det(1 - F_v t, {_\mu V_v}^{I_v})
\tag{9.8.1}
\end{equation}
est le m\^eme pour $_\lambda V$ et $_\mu V$.

Soit $v_0 \in S$, et tordons $_\lambda V$ et $_\mu V$ par un caract\`ere $\chi$ du groupe des classes d'id\`eles non ramifi\'e en $v_0$ et tr\`es ramifi\'e en les $v \in S$ autres que $v_0$ ($\chi$ est \`a valeurs dans une extension de $E$; pour l'existence, cf. 4.14).

Les hypoth\`eses faites sur $(_\lambda V, {_\mu V})$ sont stables par torsion. Pour $\chi$ convenable, apr\`es torsion, les facteurs relatifs aux $v \not= v_0$ deviennent 1. Celui relatif \`a $v_0$ subit une modification triviale, qui se lit sur $\chi(F)$.

Posons $P_v^\lambda(T) = \det(1 - F_v T, ({_\lambda V_v}^{\mathrm{ss}})^{I_v})$. Appliquant l'invariance de (9.8.1) \`a $(_\lambda V \chi, {_\mu V} \chi)$, on trouve que, pour chaque $v \in S$, on a, dans $E(T)$,
\begin{equation}
P_v^\lambda(T) / P_v^\lambda(q_v T) = P_v^\mu(T) / P_v^\mu(q_v T).
\tag{9.8.2}
\end{equation}
La transformation qui \`a une fraction rationnelle $R$ valant 1 en 0 associe $R(qT) \cdot R(T)/R(qT)$ est injective: pour les diviseurs, on a $\div(R) = \lim_{n \to \infty} \div(R(q^n T))$.

De 9.8.2, on d\'eduit donc que $P_v^\lambda(T) / P_v^\mu(T) \in E[T]$ et que $P_v^\mu(T)$ (limite simple), \`a une extension finie $K'$ pr\`es (remplacer $S$ par son image r\'eciproque). Les hypoth\`eses faites sur $_\lambda V$ et $_\mu V$ sont respect\'ees par passage de $K$ \`a $K'$ (car $\Gal(\overline{K}_v/K_v)$ est un sous-groupe ferm\'e de $\Gal(\overline{K}/K)$), on trouve que cela vaut aussi pour les repr\'esentations de $W(\overline{K}'_v/K'_v)$ d\'efinies par $_\lambda V$ et $_\mu V$. Prenant pour $K'_v$ les diverses extensions telles que les semi-simplifi\'ees de ces repr\'esentations soient non ramifi\'ee, on trouve que, pour tout $\sigma \in I_v$, on a $F\sigma \in W(\overline{K}_v/K_v)$:
\[
\Tr(\sigma, {_\lambda V_v}^{\mathrm{ss}}) = \Tr(\sigma, {_\mu V_v}^{\mathrm{ss}}).
\]
En dehors de $I_v$, les caract\`eres des repr\'esentations ${_\lambda V_v}^{\mathrm{ss}}$ et ${_\mu V_v}^{\mathrm{ss}}$ co\"ncident donc. Les m\^emes arguments, o\`u on ne retient de 9.8 que l'\'egalit\'e des degr\'es, donnent que pour tout sous-groupe ouvert $J$ de $I_v$ et tout caract\`ere $\chi$ de $J$, $\langle \chi, {_\lambda V_v}|J \rangle = \langle \chi, {_\mu V_v}|J \rangle$. D'apr\`es le th\'eor\`eme de Brauer, ceci implique que ${_\lambda V_v}^{\mathrm{ss}}$ et ${_\mu V_v}^{\mathrm{ss}}$ ont m\^eme caract\`ere. Les caract\`eres de $W(\overline{K}_v/K_v)$, et ceci ach\`eve la d\'emonstration.
\end{proof}

\begin{proposition}[9.9]
Soit $(_\lambda V)_{\lambda \in \mathcal{L}}$ un syst\`eme infini de repr\'esentations $\lambda$-adiques. On suppose que pour toute (ou presque toute: 9.8) place $v$ de $K$, les semi-simplifi\'ees des $_\lambda V_v$ sont compatibles. Alors, pour chaque $\lambda$,
\begin{equation}
Z^{\mathrm{ss}}(_\lambda V,t) = \varepsilon^{\mathrm{ss}}(_\lambda V,t) \cdot Z^{\mathrm{ss}}(_\lambda V^*, pt^{-1}).
\tag{9.9.0}
\end{equation}
\end{proposition}

Posons
\begin{align*}
Z^{\mathrm{ss}}(_\lambda V,t) &= \prod_v \det(1 - F_v t^{\deg(v)}, ({_\lambda V_v}^{\mathrm{ss}})^{I_v})^{-1}, \\
\varepsilon^{\mathrm{ss}}(_\lambda V,t) &= \prod_v \varepsilon({_\lambda V_v}^{\mathrm{ss}}, \psi_v, dx_v, t).
\end{align*}
Un calcul local montre que (9.9.0) \'equivaut \`a
\begin{equation}
Z(j_!^g({_\lambda V}^{\mathrm{ss}}),t) = \varepsilon(j_!^g({_\lambda V}^{\mathrm{ss}}),t) \cdot Z(Rj_!^g({_\lambda V}^{\mathrm{ss}*}),t^{-1}).
\tag{9.9.1}
\end{equation}
(Le quotient des facteurs locaux en $v$ des deux membres de (9.9.0) ou (9.9.1) est additif en ${_\lambda V_v}$, et ces quotients co\"ncident pour les semi-simplifi\'ees par d\'efinition de $\varepsilon^{\mathrm{ss}}$ et $Z^{\mathrm{ss}}$.)


\section{La th\'eorie de Grothendieck}

\begin{theorem}[10.1 --- Grothendieck]
Soient $X$ une courbe projective lisse et connexe sur $k = \mathbb{F}_q$, $\overline{k}$ une cl\^oture alg\'ebrique de $k$, $\overline{X} = X \otimes_k \overline{k}$, $\ell$ un nombre premier $\not= p$ et $F$ un faisceau constructible de $\mathbb{Q}_\ell$-vectoriels sur $X$. On suppose $F$ de poids $\leq n$ (i.e. pour tout point ferm\'e $x$ de $X$, les valeurs propres de $F_x$ agissant sur $F_{\bar{x}}$ sont de valeur absolue $\leq N(x)^{n/2}$). Alors
\begin{enumerate}
\item[(i)] $H^i(\overline{X}, F)$ est de poids $\leq n+i$.
\item[(ii)] $H^i_c(\overline{X}, F)$ est de poids $\leq n+i$.
\item[(iii)] $H^i(\overline{X}, F)$ est pur de poids $n+i$ si $F$ est pur de poids $n$.
\item[(iv)] Si $X$ est affine, $H^i(\overline{X}, F)$ est de poids $\leq n+i-1$ pour $i \geq 1$.
\end{enumerate}
\end{theorem}

\begin{proof}
(i) Pour $k$ fini, le th\'eor\`eme r\'esulte de ce que le polyn\^ome caract\'eristique de Frobenius agissant sur $H^i(\overline{X}, F)$ est \`a coefficients entiers (th\'eor\`eme de r\'ealisabilit\'e) et que les valeurs propres de Frobenius sont de valeur absolue $q^{(n+i)/2}$ (rappel\'e dans Weil I 6.6, prouv\'e dans SGA 7 XXI). Pour le cas g\'en\'eral, cf. [5] 3.3.1.

(ii) r\'esulte de (i) par dualit\'e de Poincar\'e.

(iii) Pour $F$ pur, (i) et la dualit\'e montrent que $H^i(\overline{X}, F)$ est de poids $\geq n+i$. D'apr\`es (i), il est aussi de poids $\leq n+i$.

(iv) est profond. Si $X$ est la compl\'ementaire dans une courbe projective lisse $\overline{X}$ d'un ensemble fini $S$ de points, (iv) r\'esulte de la suite exacte de cohomologie
\[
\cdots \to H^i_c(\overline{X} - S, F) \to H^i(\overline{X}, F) \to \bigoplus_{s \in S} H^i(s, F) \to \cdots
\]
et de ce que, pour $i > 0$, $H^i(s, F) = 0$.
\end{proof}

\begin{definition}[10.2]
Si $j: U \hookrightarrow X$ est l'inclusion d'un ouvert dense, et $F$ un faisceau sur $U$, on pose
\[
j_!^g F = \mathrm{Im}(j_! F \to j_* F).
\]
Si $F$ est irr\'eductible, on a soit $j_!^g F = j_! F$ ($F$ ramifi\'e en au moins un point de $X - U$), soit $j_!^g F = j_* F$ ($F$ non ramifi\'e en tout point de $X - U$, et $j_* F$ est un faisceau lisse sur $X$). Dans les deux cas, $j_!^g F$ est un faisceau pervers (\`a un d\'ecalage pr\`es).
\end{definition}

\begin{definition}[10.3]
Soit $F$ un faisceau constructible de $\mathbb{Q}_\ell$-vectoriels sur $X$ lisse connexe de dimension 1. On d\'efinit le conducteur de Swan global de $F$ par
\[
\sw(F) = \sum_{x \in X} \sw_x(F).
\]
On a $\sw(F) \geq 0$, et $\sw(F) = 0$ si et seulement si $F$ est mod\'er\'ement ramifi\'e. Si $j: U \hookrightarrow X$ est l'inclusion d'un ouvert dense o\`u $F$ est lisse, on a $\sw(F) = \sw(j_* j^* F)$.
\end{definition}

\begin{theorem}[10.4 --- Formule de Grothendieck-Ogg-Shafarevitch]
Soit $F$ un faisceau constructible de $\mathbb{Q}_\ell$-vectoriels sur une courbe projective lisse connexe $X$ sur $k$ alg\'ebriquement clos. On a
\[
\chi(X, F) = \rank(F) \cdot \chi(X) - \sum_{x \in X} (\rank(F) - \dim F_x + \sw_x(F)).
\]
\end{theorem}

\begin{proof}
Le th\'eor\`eme r\'esulte de la formule des traces de Lefschetz pour l'endomorphisme de Frobenius, appliqu\'ee \`a $F$ et aux faisceaux constants sur $X$, et de la formule de Grothendieck pour la caract\'eristique d'Euler-Poincar\'e.
\end{proof}

\begin{corollary}[10.5]
Sous les hypoth\`eses de 10.4, supposons de plus $F$ lisse sur $X - S$, pour $S$ fini. Alors
\[
\chi(X, j_* j^* F) = \rank(F) \cdot \chi(X - S) + \sum_{s \in S} (\dim F_s^{I_s} + \sw_s(F)).
\]
\end{corollary}

\begin{definition}[10.6]
Soit $F$ un faisceau lisse sur un ouvert $j: U \hookrightarrow X$ d'une courbe projective lisse connexe. On pose
\[
\varepsilon^{\mathrm{Gr}}(X, j_! F) = \prod_v \varepsilon_0(F_{\bar{\eta}}, \psi, dx, t) \cdot \det(-\mathrm{Frob} \cdot t^{\deg(v)}, (j_* F)_{\bar{v}})^{-1}.
\]
On a $\varepsilon^{\mathrm{Gr}} \in \overline{\mathbb{Q}}_\ell[t, t^{-1}]^*$.
\end{definition}

\begin{theorem}[10.7 --- Grothendieck]
Avec les notations pr\'ec\'edentes,
\[
Z(U, F, t) = \varepsilon^{\mathrm{Gr}}(X, j_! F) \cdot Z(U, F^*, p^{-1} t^{-1}).
\]
\end{theorem}

\section{Appendice}

\subsection*{A. Sommes de Gauss}

Soient $k$ un corps fini \`a $q$ \'el\'ements, $\psi$ un caract\`ere additif non trivial de $k$ et $\chi$ un caract\`ere multiplicatif de $k^*$. On pose
\[
\tau(\chi, \psi) = -\sum_{x \in k^*} \chi^{-1}(x) \psi(x).
\]

\begin{proposition}[A.1]
Si $\chi$ est trivial, $\tau(\chi, \psi) = 1$. Si $\chi$ est non trivial, $|\tau(\chi, \psi)| = \sqrt{q}$.
\end{proposition}

\begin{proposition}[A.2 --- Hasse-Davenport]
Soient $k'$ une extension de degr\'e $n$ de $k$ et $\chi$ un caract\`ere de $k^*$. On a
\[
\tau(\chi \circ N_{k'/k}, \psi \circ \Tr_{k'/k}) = (-1)^{n-1} \tau(\chi, \psi)^n.
\]
\end{proposition}

\subsection*{B. Repr\'esentations de Swan}

Soient $K$ un corps local non archim\'edien, $I$ le groupe d'inertie et $P$ le sous-groupe de wild inertia. Pour $G$ un quotient fini de $I$, on d\'efinit la repr\'esentation de Swan $\sw_G$ comme la repr\'esentation virtuelle
\[
\sw_G = \sum_{i \geq 1} \frac{|G_i|}{|G|} \cdot \Ind_{G_i}^G(\chi_i)
\]
o\`u $G_i$ sont les sous-groupes de ramification en num\'erotation sup\'erieure et $\chi_i$ sont des caract\`eres fid\`eles de $G_i/G_{i+1}$.

\begin{proposition}[B.1]
$\sw_G$ est une repr\'esentation virtuelle de $G$ \`a valeurs enti\`eres. Pour tout caract\`ere $\chi$ de $G$, on a $\sw(\chi) = \langle \sw_G, \chi \rangle$.
\end{proposition}

\begin{thebibliography}{20}

\bibitem{1} E. Artin, J. Tate, \textit{Class Field Theory}, Benjamin, New York, 1967.

\bibitem{2} P. Deligne, Les constantes des \'equations fonctionnelles des fonctions $L$, \textit{Modular functions of one variable II}, Lecture Notes in Math. 349, Springer-Verlag, 1973, pp. 55--68.

\bibitem{3} P. Deligne, La conjecture de Weil I, \textit{Publ. Math. IHES} 43 (1974), 273--307.

\bibitem{4} S. Gelbart, H. Jacquet, Forms of $\GL(2)$ from the analytic point of view, \textit{Proc. Sympos. Pure Math.} 33 (1979), 213--251.

\bibitem{5} N. Katz, W. Messing, Some consequences of the Riemann hypothesis for varieties over finite fields, \textit{Invent. Math.} 23 (1974), 73--77.

\bibitem{6} R. Langlands, On the functional equation of the Artin $L$-functions, \textit{Yale University preprint} (unpublished).

\bibitem{7} R. Langlands, On the classification of irreducible representations of real algebraic groups, \textit{Mimeographed notes}, Institute for Advanced Study, 1973.

\bibitem{8} H. Jacquet, R. Langlands, \textit{Automorphic forms on $\GL(2)$}, Lecture Notes in Math. 114, Springer-Verlag, 1970.

\bibitem{9} R. Langlands, \textit{On the functional equations satisfied by Eisenstein series}, Lecture Notes in Math. 544, Springer-Verlag, 1976.

\bibitem{10} J.-P. Serre, Repr\'esentations lin\'eaires des groupes finis, Hermann, Paris, 1967.

\bibitem{11} J.-P. Serre, Facteurs locaux des fonctions z\^eta des vari\'et\'es alg\'ebriques (d\'efinitions et conjectures), \textit{S\'em. Delange-Pisot-Poitou} 1969/70, exp. 19.

\bibitem{12} J.-P. Serre, J. Tate, Good reduction of abelian varieties, \textit{Ann. of Math.} 88 (1968), 492--517.

\bibitem{13} J.-P. Serre, \textit{Abelian $\ell$-adic representations and elliptic curves}, Benjamin, New York, 1968.

\bibitem{14} J.-P. Serre, Corps locaux, Hermann, Paris, 1962.

\bibitem{15} A. Weil, \textit{Basic Number Theory}, Springer-Verlag, 1967.

\bibitem{16} A. Weil, Dirichlet series and automorphic forms, \textit{Lecture Notes in Math.} 189, Springer-Verlag, 1971.

\bibitem{17} A. Weil, Numbers of solutions of equations in finite fields, \textit{Bull. Amer. Math. Soc.} 55 (1949), 497--508.

\bibitem{18} A. Weil, Jacobi sums as ``Gr\"ossencharaktere'', \textit{Trans. Amer. Math. Soc.} 73 (1952), 487--495.

\bibitem{19} H. Weyl, \textit{The Classical Groups}, Princeton University Press, 1946.

\bibitem{20} A. Grothendieck, Formule de Lefschetz et rationalit\'e des fonctions $L$, \textit{S\'em. Bourbaki} 1964/65, exp. 279.

\bibitem{21} P. Deligne, Le formalisme des cycles \'evanescents, SGA 7 XIII, \textit{Lecture Notes in Math.} 340, Springer-Verlag, 1973.

\bibitem{22} P. Deligne, La formule de dualit\'e globale, SGA 4 XVIII, \textit{Lecture Notes in Math.} 305, Springer-Verlag, 1973.

\bibitem{23} N. Katz, An overview of Deligne's proof of the Riemann hypothesis for varieties over finite fields, \textit{Proc. Sympos. Pure Math.} 28 (1976), 275--305.

\bibitem{24} G. Laumon, Transformation de Fourier, constantes d'\'equations fonctionnelles et conjecture de Weil, \textit{Publ. Math. IHES} 65 (1987), 131--210.

\bibitem{25} P. Deligne, Le lemme de Gabber, SGA 4 1/2, \textit{Lecture Notes in Math.} (Ast\'erisque).

\bibitem{26} J.-L. Brylinski, Th\'eor\`eme de r\'eciprocit\'e pour les suites spectrales, \textit{Th\`ese de 3\`eme cycle}, Universit\'e Paris VII, 1976.

\bibitem{27} P. Deligne, Th\'eor\`eme de Lefschetz et crit\`eres de d\'eg\'en\'erescence de suites spectrales, \textit{Publ. Math. IHES} 35 (1968), 107--126.

\bibitem{28} T. Saito, The sign of the functional equation of the $L$-function of an orthogonal motive, \textit{Invent. Math.} 120 (1995), 119--142.

\bibitem{29} F. Shahidi, Local coefficients as Artin factors for real groups, \textit{Duke Math. J.} 52 (1985), 973--1007.

\bibitem{30} J. Tate, Number theoretic background, \textit{Proc. Sympos. Pure Math.} 33 (1979), 3--26.

\end{thebibliography}

\end{document}
